\section{Numerical experiments}\label{sec:numerics}

\subsection{Models, methods and measurement}\label{sec:setup}
\paragraph*{Models.} Two families of leaky diffusion networks are used, both of the form $\dot\bx=-\mathcal{C}^{-1}(L+G)\bx+B\bu$ with $L$ a weighted graph Laplacian, $G=\mathrm{diag}(\ell_j)$ a leakage to ground and $\mathcal{C}=\mathrm{diag}(c_j)$ capacities, so that $d_j=-(\sum_kg_{jk}+\ell_j)/c_j$ and $R_{jk}=g_{jk}/c_j$. The comparison ratio of state $j$ is $\sum_kg_{jk}/(\sum_kg_{jk}+\ell_j)$, independent of $c_j$.
(a) \emph{Thermal network}: an $L\times L$ grid of lumped thermal masses with conductances $g_{jk}\in[0.75,1.25]$ between nearest neighbors, convective losses $\ell_j\in\ell_0[0.75,1.25]$ and capacities $c_j\in[0.75,1.25]$ (or log-uniform over $s_C$ decades in Sec.~\ref{sec:stiff}), all drawn from a hash of the indices. This is the standard RC model of in-plane heat conduction. The distance oracle is the Manhattan distance.
(b) \emph{Implicit sparse random network}: $N=2^n$ states, $n\le64$, coupled through $K$ pseudorandom bijections $\pi_k$ of $\{0,\dots,N-1\}$ realized as four-round Feistel networks, so that state $j$ is linked to $\pi_k(j)$ and $\pi_k^{-1}(j)$ (degree $2K$). Conductances, losses and capacities are drawn as in (a). Such graphs are locally tree-like up to distances of order $\log N$, as sparse random networks are, and the dynamics is the grounded-Laplacian (leader--follower) network model of Refs.~\cite{mesbahi2010,pasqualetti2014}. We use it as a structural stand-in for large sparse interaction networks and do not attach it to a particular physical system. Nothing about it is stored; every oracle call evaluates $O(K)$ Feistel rounds.
The sensor and the actuator are collocated at one state (a single heater with a thermocouple, a driver node that is also observed) unless stated otherwise. Outputs are the impulse kernel and the step response on $t_k=0.05k$, $k\le100$.

\paragraph*{Methods.} The walk sum is the C++ implementation of Algorithm~1. Global comparators run on the explicitly exported matrix: the scaled Taylor method of Al-Mohy and Higham as implemented in SciPy's \texttt{expm\_multiply}~\cite{almohy2011,virtanen2020}, advanced one output interval at a time so that only a few state vectors are live (the step response uses the augmented generator); an Arnoldi exponential with $m=30$ and step-size control on the usual a posteriori estimate~\cite{hochbruck1997,sidje1998}; and, for the stiff family, SciPy's variable-order BDF integrator with the exact sparse Jacobian~\cite{shampine1997}. Global methods are run forward from the actuator; the adjoint propagation with $A^{\mathsf T}$ from the sensor has the same cost for one sensor and one actuator, and on the $101\times101$ grid with sensor and actuator three sites apart the two agreed to $7\times10^{-20}$ on outputs of order $10^{-4}$. Two local comparators carry the same guarantee as the walk sum (Sec.~\ref{sec:local}). The \emph{certified ball} uses the smallest radius whose certificate meets the target; the induced principal submatrix is extracted by breadth-first search. The \emph{adaptive set} grows a retained set by thresholded expansion until the certificate of Proposition~\ref{prop:local} is below 0.9 times the target. Both are propagated by the compiled Taylor integrator of Sec.~\ref{sec:local}, in the same binary and with the same compiler as the walk sum, with the remaining part of the target as the certified time-truncation budget. Where a smaller ball is solved with SciPy instead, this is stated. All three methods compute the impulse kernel and the step response, and the accuracy target applies to the step-response certificate.

\paragraph*{Distance information.} Unless stated otherwise, the walk sum and the adaptive set use the same distance lower bounds: the Manhattan distance on the grid, and on the random networks a breadth-first map of radius $r_d=3$ around the actuator, built at start-up and counted in each process's time and memory. The ball's certificate uses the exact radius of its own breadth-first search. Section~\ref{sec:crossover} also reports the walk sum without any distance map, which is its most memory-lean configuration.

\paragraph*{Time to a certificate.} The walk sum does not know in advance which pruning threshold meets a target. We lower the threshold until the certificate meets it, starting each target from the threshold that met the previous, looser one, and we report two times. The first is the traversal at the final threshold, which is what a user who knows the threshold pays. The second is the cumulative time of all attempts for that model up to and including the target, including the construction of the distance map in every attempt, which is what a user who starts from scratch pays. The local methods are timed from scratch: set construction, certificate computation and time integration.

\paragraph*{Measurement.} Every method runs in its own process. Peak memory is the high-water resident set size of that process (\texttt{VmHWM}, which is reset at \texttt{exec}); for Python processes we also report the high-water mark after importing NumPy and SciPy and quote the difference as incremental memory. Times are wall-clock times of the computation phase, with model construction, export and loading reported separately. All computations are single-threaded (numerical libraries pinned to one thread) on one core of an Intel Xeon at 2.8~GHz with 7.7~GB of memory, using g++~13.3 with \texttt{-O3 -march=native}, Python~3.13, NumPy~2.5 and SciPy~1.18. Accuracy comparisons are made at matched certificates, the only accuracy a user can rely on without a reference; the observed errors are reported as well. Runs that were not performed, for instance state-vector methods at $N=2^{64}$, are marked as such and never extrapolated. Times are single measurements unless stated otherwise. On this virtual machine identical runs repeated on an otherwise idle system differed by up to about 40\% for sub-second runs and by a few percent for runs of several seconds, and an earlier pass of the comparison in Table~\ref{tab:crossover}, taken while short unrelated jobs ran on the second core, was up to twice slower for some local runs; the rows of that table with non-negligible times were therefore re-timed three times on an idle system and are reported as medians. Differences of less than a factor of two between single measurements should not be read as significant.
